%% Drafted from the mapped corpus by scripts/draft_topics.py.
% Model: gpt-5.6-terra; source characters: 21362.
\begin{konspekt}
\kselect{На изпита ще се падне \textbf{или} алгоритъмът на Прим, \textbf{или} алгоритъмът на Крускал.}
\kreq{Дефиниция на задачата}{какъв вид графи се разглеждат и какво се търси --- \kref{7.1}.}
\kreq{Теорема с доказателство}{за съгласуваното множество --- условия за нарастване на подмножество на МПД --- \kref{7.2}.}
\kreq{Псевдокод, коректност, сложност при различни структури от данни}{алгоритъм на Прим --- \kref{7.3}.}
\kreq{Псевдокод и коректност}{алгоритъм на Крускал --- \kref{7.4}.}
\kreq{Структура от данни}{Union--Find и приложението му за ефикасна реализация --- \kref{7.5}.}
\kreq{Сложност по време}{на алгоритъма на Крускал --- \kref{7.6}.}
\kreq{Литература}{[31].}
\end{konspekt}

\section{Задача за минимално покриващо дърво}\label{s:7.1}

\begin{defn}
Тегловен граф е наредената двойка $(G,\omega)$, където $G=(V,E)$ е граф, а
\[
\omega:E\to\mathbb{R}
\]
е теглова функция, която задава тегло, цена или дължина на всяко ребро.
\end{defn}

В задачата за минимално покриващо дърво се разглежда неориентиран свързан тегловен граф. Ориентацията не е част от задачата, защото понятието „дърво“ тук описва свързан ацикличен неориентиран подграф.

\begin{defn}
Нека $G=(V,E)$ е граф. Покриващ подграф на $G$ е всеки подграф от вида
\[
G'=(V,E'), \qquad E'\subseteq E.
\]
Покриващо дърво на $G$ е покриващ подграф, който е дърво.
\end{defn}

Тъй като графът $G$ е свързан, той има поне едно покриващо дърво. Всяко покриващо дърво съдържа всички върхове на $G$ и точно $|V|-1$ ребра.

\begin{defn}
Нека $G=(V,E)$ е свързан тегловен граф с теглова функция $\omega$, а $T$ е покриващо дърво на $G$. Тегло на $T$ наричаме сумата
\[
\omega(T)=\sum_{e\in E(T)}\omega(e).
\]
Покриващото дърво $T$ е \emph{минимално покриващо дърво}, съкратено МПД, ако за всяко покриващо дърво $T'$ на $G$ е изпълнено
\[
\omega(T)\leq \omega(T').
\]
\end{defn}

\begin{keydefn}
Задачата за минимално покриващо дърво има следния вид.

Екземплярът е наредена двойка $\langle G,\omega\rangle$, където $G=(V,E)$ е неориентиран свързан граф, а $\omega:E\to\mathbb{R}$ е теглова функция.

Решението е минимално покриващо дърво $T$ на $G$.
\end{keydefn}

Теглата не е необходимо да са положителни и не е необходимо да са различни. При равни тегла могат да съществуват различни минимални покриващи дървета.

\section{Съгласувани множества и безопасни ребра}\label{s:7.2}

Алчните алгоритми за МПД изграждат множество от ребра $A$, което постепенно се разширява. Основният въпрос е кога към $A$ може безопасно да се добави ново ребро.

\begin{defn}
Нека $G=(V,E)$ е тегловен граф и $A\subseteq E$. Казваме, че $A$ е \emph{съгласувано} множество, ако съществува минимално покриващо дърво $T$ на $G$, за което
\[
A\subseteq E(T).
\]
\end{defn}

\begin{defn}
Нека $A\subseteq E$ е съгласувано множество и $e\in E\setminus A$. Реброто $e$ е \emph{сигурно} за $A$, ако съществува минимално покриващо дърво $T'$ на $G$, за което
\[
A\cup\{e\}\subseteq E(T').
\]
\end{defn}

\begin{defn}
Срез в графа $G=(V,E)$ е разделяне на множеството от върхове на две непразни части:
\[
S=\{V_1,V_2\},\qquad V_1\cap V_2=\varnothing,\qquad V_1\cup V_2=V.
\]
Ребро $e=\{u,v\}$ \emph{пресича} среза $S$, ако единият му край е във $V_1$, а другият --- във $V_2$.

Срезът $S$ е \emph{съобразен} с множество $A\subseteq E$, ако нито едно ребро от $A$ не пресича $S$.

Ребро, пресичащо среза, е \emph{леко} за този срез, ако теглото му е минимално измежду теглата на всички ребра, които пресичат среза.
\end{defn}

\begin{keythm}[Теорема за съгласуваното множество]
Нека $G=(V,E)$ е свързан тегловен граф с теглова функция $\omega$. Нека $A\subseteq E$ е съгласувано множество и $S$ е срез, съобразен с $A$. Ако $e$ е произволно леко ребро, което пресича $S$, то $e$ е сигурно за $A$.
\end{keythm}

\begin{proof}
Тъй като $A$ е съгласувано, съществува минимално покриващо дърво $T$, за което
\[
A\subseteq E(T).
\]
Ако $e\in E(T)$, твърдението следва веднага.

Нека сега $e\notin E(T)$. Понеже $T$ е дърво, между двата края на $e$ има точно един път в $T$. След добавяне на $e$ към $T$ се получава точно един цикъл $C$.

Реброто $e$ пресича среза $S$. Следователно, при обхождането на цикъла $C$ трябва да се премине поне още веднъж от едната част на среза в другата. Значи в $C$ съществува ребро
\[
e'\neq e,
\]
което също пресича $S$.

Тъй като срезът е съобразен с $A$, реброто $e'$ не принадлежи на $A$. Премахваме $e'$ от цикъла $C$ и получаваме графа
\[
T'=T-\{e'\}+\{e\}.
\]
Това е покриващо дърво: след премахването на едно ребро от единствения цикъл графът остава свързан и става ацикличен. Освен това
\[
A\cup\{e\}\subseteq E(T'),
\]
защото $e'\notin A$.

Тъй като $e$ е леко ребро за среза, а $e'$ също пресича този срез, имаме
\[
\omega(e)\leq \omega(e').
\]
Следователно
\[
\omega(T')
 =\omega(T)-\omega(e')+\omega(e)
 \leq \omega(T).
\]
Но $T$ е минимално покриващо дърво, следователно $\omega(T)\leq\omega(T')$. Оттук
\[
\omega(T')=\omega(T),
\]
така че $T'$ също е минимално покриващо дърво. Значи $e$ е сигурно за $A$.
\end{proof}

Теоремата обосновава общия алчен подход:

\begin{verbatim}
Generic_MST(G = (V,E), omega)
A <- empty set                 // избраните ребра; винаги гора
while |A| < |V| - 1 do             // "A още не е покриващо дърво"
    избери ребро e, сигурно за A
    A <- A union {e}
return A
\end{verbatim}

Условието на цикъла е записано чрез броя ребра, защото гора върху $|V|$ върха е покриващо дърво тогава и само тогава, когато има точно $|V|-1$ ребра.

Ако на всяка стъпка $A$ е съгласувано и се добавя сигурно ребро, то след добавянето множеството остава съгласувано. Когато то вече образува покриващо дърво, това дърво трябва да е минимално.

Схемата не казва как се намира сигурно ребро --- това е единствената разлика между двата алгоритъма по-долу. Прим държи един срез между построеното дърво и останалите върхове и взема леко ребро за него; Крускал разглежда ребрата в ненамаляващ ред и за всяко прието ребро подходящият срез се намира впоследствие.

\section{Алгоритъм на Прим}\label{s:7.3}

Алгоритъмът на Прим строи едно дърво, започвайки от избран начален връх. На всяка стъпка той добавя най-лекото налично ребро, което свързва вече построената част на дървото с връх извън нея.

\subsection{Идея и поддържани данни}\label{s:7.3.1}

Нека $r\in V$ е начален връх. Поддържа се приоритетна опашка $Q$ от върховете, които все още не са добавени в дървото. За всеки връх $v$ се пазят две полета:

\begin{itemize}
\item $v.\mathit{key}$ --- теглото на най-лекото намерено досега ребро, което свързва $v$ с вече добавен връх; това е приоритетът на $v$ в опашката;
\item $v.\pi$ --- другият край на същото това ребро, тоест предшественикът на $v$ в строящото се дърво.
\end{itemize}

Ако още не е намерено ребро от дървото към $v$, тогава $v.\mathit{key}=\infty$ и $v.\pi=\textsc{Nil}$. За началния връх се поставя ключ $0$, за да бъде избран първи.

Дървото никъде не се пази явно. То се пази разпределено, в указателите $\pi$: щом връх бъде извлечен от опашката, реброто $\{v,v.\pi\}$ е окончателно избрано и вече не се променя.

\begin{keydefn}
Нека $X=V\setminus Q$ е множеството от вече извлечените от опашката върхове. Връх от $Q$ е граничен, ако има съсед в $X$; останалите върхове от $Q$ са неизвестни. Граничните върхове имат краен ключ, а неизвестните --- ключ $\infty$.
\end{keydefn}

Върховете от $X$ са точно тези, за които реброто към дървото е избрано окончателно; $\textsc{Extract-Min}$ винаги избира граничен връх, докато $Q$ съдържа поне един такъв, тъй като графът е свързан.

\subsection{Използвани операции}\label{s:7.3.2}

Псевдокодът извиква операции на приоритетна опашка. Те не са част от алгоритъма на Прим: всяка реализация, поддържаща ги, е допустима, а изборът на реализация влияе само на сложността.

\begin{defn}[Приоритетна опашка по минимум]
Приоритетна опашка е структура от данни, която пази множество от елементи, всеки с числов ключ, и поддържа операциите:

\begin{itemize}
\item \verb|Insert(Q,x)| --- добавя елемента $x$ с текущия му ключ;
\item \verb|Extract-Min(Q)| --- премахва от $Q$ и връща елемент с минимален ключ;
\item \verb|Decrease-Key(Q,x,k)| --- заменя ключа на $x$ с новата стойност $k$, която не е по-голяма от старата, и възстановява вътрешния ред на структурата.
\end{itemize}
\end{defn}

В псевдокода се използват още две означения:

\begin{itemize}
\item $Q\leftarrow V$ --- построяване на опашка от всички върхове с текущите им ключове, тоест $|V|$ операции \verb|Insert| или еднократно построяване;
\item $\mathit{Adj}[u]$ --- списъкът на съседство на $u$, тоест списъкът на върховете, съседни на $u$.
\end{itemize}

\subsection{Псевдокод}\label{s:7.3.3}

\begin{verbatim}
MST_Prim(G = (V,E), omega, r)
// 1. никой връх още не е свързан с дървото
for each v in V do
    v.key <- infinity        // най-лекото ребро от дървото до v
    v.pi  <- NIL             // другият край на това ребро
r.key <- 0                   // за да бъде r извлечен пръв
Q <- V                       // всички върхове чакат в опашката
// 2. всяка итерация добавя по един връх към дървото
while Q is not empty do
    u <- Extract-Min(Q)      // най-евтино присъединимият връх
    for each v in Adj[u] do  // ребрата от u към съседите му
        if v in Q and omega({u,v}) < v.key then
            v.pi  <- u       // u е новият най-евтин вход към v
            v.key <- omega({u,v})
            Decrease-Key(Q, v, v.key)
// 3. дървото се чете от указателите pi
return {{v, v.pi} : v in V and v.pi != NIL}
\end{verbatim}

Означенията в псевдокода са следните.

\begin{itemize}
\item $G=(V,E)$, $\omega$ и $r$ са входът: свързан неориентиран граф, теглова функция и начален връх.
\item $Q$ е приоритетната опашка на върховете, които още не са в дървото; $V\setminus Q$ са вече добавените върхове.
\item Тялото на \verb|while| се изпълнява по веднъж за всеки връх, защото \verb|Extract-Min| премахва връх и нищо не се връща обратно в $Q$.
\item Върнатото множество съдържа по едно ребро $\{v,v.\pi\}$ за всеки връх, различен от $r$; само за $r$ остава $v.\pi=\textsc{Nil}$.
\end{itemize}

В неориентиран граф всяко ребро се среща в списъците на съседство на двата си края. Проверката $v\in Q$ може да се реализира за константно време с отделна булева или битова информация дали върхът е в опашката.

\subsection{Пример}\label{s:7.3.4}

\begin{example}
Разглеждаме графа от фигура~\ref{fig:mst-example}, започвайки от връх $a$.

\begin{lecfigure}[Примерен тегловен граф. Удебелените ребра образуват минималното покриващо дърво с тегло $13$.\label{fig:mst-example}]
\begin{tikzpicture}[scale=1.25,
  vtx/.style={circle, draw=defframe, fill=defback, inner sep=0pt,
              minimum size=6.5mm, font=\small},
  eplain/.style={gray!65, line width=0.5pt},
  etree/.style={thmframe, line width=1.6pt},
  wlab/.style={font=\footnotesize, inner sep=1.4pt, fill=white}]
  \node[vtx] (a) at (0,1)  {$a$};
  \node[vtx] (b) at (1.9,2){$b$};
  \node[vtx] (c) at (1.9,0){$c$};
  \node[vtx] (d) at (3.8,2){$d$};
  \node[vtx] (e) at (3.8,0){$e$};
  \node[vtx] (f) at (5.7,1){$f$};
  \draw[eplain] (a) -- node[wlab] {4} (b);
  \draw[etree]  (a) -- node[wlab] {3} (c);
  \draw[etree]  (b) -- node[wlab] {1} (c);
  \draw[etree]  (b) -- node[wlab] {2} (d);
  \draw[eplain] (c) -- node[wlab] {4} (d);
  \draw[etree]  (c) -- node[wlab] {5} (e);
  \draw[eplain] (d) -- node[wlab] {7} (e);
  \draw[eplain] (d) -- node[wlab] {6} (f);
  \draw[etree]  (e) -- node[wlab] {2} (f);
\end{tikzpicture}
\end{lecfigure}

След инициализацията $a.\mathit{key}=0$, а всички останали ключове са $\infty$. Всеки ред от таблицата отговаря на една итерация на цикъла \verb|while|: първо е извлеченият връх и реброто, което той внася в дървото, а после ключовете, с които останалите върхове чакат в $Q$.

\begin{center}
\begin{tabular}{llll}
\toprule
Извлечен $u$ & Ребро $\{u,u.\pi\}$ & Тегло & Ключове в $Q$ след обхождането \\
\midrule
$a$ & няма & & $b{:}4(a)$, $c{:}3(a)$, $d{:}\infty$, $e{:}\infty$, $f{:}\infty$ \\
$c$ & $\{a,c\}$ & $3$ & $b{:}1(c)$, $d{:}4(c)$, $e{:}5(c)$, $f{:}\infty$ \\
$b$ & $\{c,b\}$ & $1$ & $d{:}2(b)$, $e{:}5(c)$, $f{:}\infty$ \\
$d$ & $\{b,d\}$ & $2$ & $e{:}5(c)$, $f{:}6(d)$ \\
$e$ & $\{c,e\}$ & $5$ & $f{:}2(e)$ \\
$f$ & $\{e,f\}$ & $2$ & $Q=\varnothing$ \\
\bottomrule
\end{tabular}
\end{center}

Забележете двете места, на които \verb|Decrease-Key| променя вече намерено ребро: ключът на $b$ пада от $4(a)$ на $1(c)$ и ключът на $f$ пада от $6(d)$ на $2(e)$. Реброто $\{d,e\}$ с тегло $7$ никога не подобрява ключа на $e$, а $\{a,b\}$ и $\{c,d\}$ се губят при пресмятанията на съответните ключове.

Полученото дърво е
\[
\{a,c\},\ \{c,b\},\ \{b,d\},\ \{c,e\},\ \{e,f\},
\]
с общо тегло $3+1+2+5+2=13$.
\end{example}

\subsection{Коректност}\label{s:7.3.5}

\begin{keythm}
За всеки свързан претеглен неориентиран граф алгоритъмът на Прим връща минимално покриващо дърво.
\end{keythm}

\begin{proof}
Нека $A$ е множеството от ребрата, определени чрез указателите $\pi$ за върховете, извлечени преди текущата итерация, без началния връх. Ще покажем по индукция по итерациите на цикъла, че $A$ е съгласувано множество.

Преди първата итерация $A=\varnothing$. Празното множество е съгласувано, защото свързаният претеглен неориентиран граф има минимално покриващо дърво.

Да предположим, че преди дадена итерация $A$ е съгласувано, и нека $u$ е върхът, избран с \verb|Extract-Min|. Ако $u=r$, не се добавя ребро, $A$ остава празно и инвариантът се запазва. В останалите итерации $u\neq r$. Нека преди извличането на $u$
\[
X=V\setminus Q
\]
е непразното множество от върхове, вече добавени към построяваното дърво. Реброто
\[
e=\{u,u.\pi\}
\]
свързва $u\in V\setminus X$ с връх от $X$.

Разглеждаме среза
\[
S=\{X,V\setminus X\}.
\]
Той е съобразен с $A$, защото всички ребра от $A$ имат и двата си края в $X$. По време на обхожданията на списъците на съседство алгоритъмът поддържа за всеки връх $v\in Q$ теглото на най-лекото ребро между $v$ и $X$. Следователно изборът на $u$ с минимален ключ означава, че $\{u,u.\pi\}$ пресича среза $S$ и е леко ребро за него.

По теоремата за съгласуваното множество реброто $e$ е сигурно за $A$. Следователно
\[
A\cup\{e\}
\]
също е съгласувано. Индуктивната хипотеза се запазва.

След завършване на алгоритъма всички върхове са извлечени. Всеки връх, различен от $r$, има точно един предшественик, следователно върнатото множество съдържа $|V|-1$ ребра. Всяко ново ребро свързва нов връх с вече построената част, затова не може да образува цикъл и всички върхове са свързани. Следователно резултатът е покриващо дърво. Понеже множеството от неговите ребра е съгласувано, то се съдържа в някое минимално покриващо дърво; а и двете дървета имат по $|V|-1$ ребра. Значи върнатото дърво е минимално покриващо.
\end{proof}

\subsection{Сложност на алгоритъма на Прим}\label{s:7.3.6}

Нека $n=|V|$ и $m=|E|$. При списъци на съседство всички обхождания на редовете с \verb|for each| разглеждат общо $O(m)$ съседства. Всеки връх се извлича от приоритетната опашка точно веднъж. Операцията \verb|Decrease-Key| се извиква най-много по веднъж за всяко разгледано съседство, следователно най-много $O(m)$ пъти.

\begin{prop}
При реализация на $Q$ с двоична пирамида алгоритъмът на Прим работи за време
\[
O((n+m)\log n)=O(m\log n)
\]
за свързан граф.
\end{prop}

\begin{proof}
Построяването на двоичната пирамида от всички върхове отнема $O(n)$ време. Има $n$ операции \verb|Extract-Min|, всяка за $O(\log n)$ време, общо $O(n\log n)$. Операцията \verb|Decrease-Key| също отнема $O(\log n)$ и се извиква най-много $O(m)$ пъти, което дава $O(m\log n)$.

Следователно общото време е
\[
O(n+m+n\log n+m\log n)=O((n+m)\log n).
\]
Понеже при свързан граф $m\geq n-1$, получаваме $O(m\log n)$.
\end{proof}

\begin{prop}
При реализация на $Q$ с пирамида на Фибоначи алгоритъмът на Прим работи за амортизирано време
\[
O(m+n\log n).
\]
\end{prop}

\begin{proof}
При пирамида на Фибоначи операцията \verb|Decrease-Key| е с амортизирана цена $O(1)$, а \verb|Extract-Min| --- с амортизирана цена $O(\log n)$. Затова $O(m)$ намалявания на ключа струват $O(m)$, а $n$ извличания на минимум струват $O(n\log n)$.
\end{proof}

\begin{remark}
При двоична пирамида често се записва и оценката $O(m\log m)$, когато логаритъмът се изразява чрез броя на ребрата. За прост неориентиран граф $m< n^2$, откъдето $\log m=O(\log n)$; така двете оценки са съвместими в рамките на обичайните предположения за графа.
\end{remark}

\section{Алгоритъм на Крускал}\label{s:7.4}

Алгоритъмът на Крускал строи не едно растящо дърво, а растяща гора. Ребрата се разглеждат в ненамаляващ ред по тегло. Ребро се добавя точно когато краищата му принадлежат на различни дървета в текущата гора.

\subsection{Използвани операции}\label{s:7.4.1}

Псевдокодът извиква три операции върху разбиване на множеството от върхове на непресичащи се дялове. Всеки дял е една компонента на свързаност на текущата гора $(V,A)$.

\begin{itemize}
\item \verb|Make-Set(v)| --- създава нов дял, съдържащ само $v$;
\item \verb|Find-Set(v)| --- връща представителя (лидера) на дяла, съдържащ $v$; два върха са в един и същ дял тогава и само тогава, когато \verb|Find-Set| връща един и същ представител за тях;
\item \verb|Union(u,v)| --- слива дяловете на $u$ и $v$ в един.
\end{itemize}

Заедно те образуват структурата Union--Find, разгледана подробно след доказателството за коректност. Засега е достатъчна само тази спецификация: нейната реализация не влияе на коректността на Крускал, а само на сложността.

\subsection{Псевдокод}\label{s:7.4.2}

\begin{verbatim}
MST_Kruskal(G = (V,E), omega)
A <- empty set                   // избраните ребра
for each v in V do
    Make-Set(v)                  // всеки връх е отделен дял
sort E in nondecreasing order by omega(e)
for each edge {u,v} in E do      // в този сортиран ред
    if Find-Set(u) != Find-Set(v) then    // различни компоненти
        A <- A union {{u,v}}     // не затваря цикъл: приема се
        Union(u,v)               // компонентите се сливат
return A
\end{verbatim}

Проверката \verb|Find-Set(u) != Find-Set(v)| е точно проверката дали реброто $\{u,v\}$ би затворило цикъл: цикъл се получава тогава и само тогава, когато двата края вече са свързани чрез приетите ребра. Затова липсва отделна проверка за цикъл --- дяловете я вършат.

При свързан граф цикълът приема точно $|V|-1$ ребра, така че изпълнението може да спре веднага щом $|A|=|V|-1$; останалите ребра неизбежно биха затворили цикъл.

\subsection{Пример}\label{s:7.4.3}

\begin{example}
За графа от фигура~\ref{fig:mst-example} сортираният ред на ребрата е
\[
\{b,c\}{:}1,\ \{b,d\}{:}2,\ \{e,f\}{:}2,\ \{a,c\}{:}3,\ \{a,b\}{:}4,\ \{c,d\}{:}4,\ \{c,e\}{:}5,\ \{d,f\}{:}6,\ \{d,e\}{:}7.
\]

\begin{center}
\begin{tabular}{llll}
\toprule
Ребро & Тегло & Действие & Компоненти след стъпката \\
\midrule
$\{b,c\}$ & $1$ & приема се & $\{a\}$, $\{b,c\}$, $\{d\}$, $\{e\}$, $\{f\}$ \\
$\{b,d\}$ & $2$ & приема се & $\{a\}$, $\{b,c,d\}$, $\{e\}$, $\{f\}$ \\
$\{e,f\}$ & $2$ & приема се & $\{a\}$, $\{b,c,d\}$, $\{e,f\}$ \\
$\{a,c\}$ & $3$ & приема се & $\{a,b,c,d\}$, $\{e,f\}$ \\
$\{a,b\}$ & $4$ & отхвърля се & без промяна \\
$\{c,d\}$ & $4$ & отхвърля се & без промяна \\
$\{c,e\}$ & $5$ & приема се & $\{a,b,c,d,e,f\}$ \\
\bottomrule
\end{tabular}
\end{center}

Приети са $5=|V|-1$ ребра, така че ребрата $\{d,f\}$ и $\{d,e\}$ вече не се разглеждат. Полученото дърво съвпада с намереното от Прим и има тегло $13$, макар че двата алгоритъма добавят ребрата в различен ред.
\end{example}

\subsection{Коректност}\label{s:7.4.4}

\begin{keythm}
Алгоритъмът на Крускал връща минимално покриващо дърво.
\end{keythm}

\begin{proof}
В началото $A=\varnothing$, следователно $A$ е съгласувано. Нека разгледаме стъпка, в която алгоритъмът приема реброто
\[
e=\{u,v\}.
\]
Условието
\[
\operatorname{Find\mbox{-}Set}(u)\neq\operatorname{Find\mbox{-}Set}(v)
\]
означава, че $u$ и $v$ са в различни компоненти на гората $(V,A)$. Нека $C$ е компонентата, съдържаща $u$, и разгледаме среза
\[
S=\{C,V\setminus C\}.
\]
Този срез е съобразен с $A$, тъй като по определение няма ребро от $A$, което да свързва различни компоненти на гората.

Реброто $e$ пресича среза. То е леко за този срез: ако имаше пресичащо ребро с по-малко тегло, то би се намирало по-рано в сортирания ред. Когато то е било разгледано, краищата му са били в различни компоненти, защото компонентите могат само да се сливат, а не да се разделят. Следователно това по-леко ребро би било добавено в $A$, което противоречи на съобразеността на среза с $A$.

От теоремата за съгласуваното множество следва, че $e$ е сигурно за $A$. Следователно след всяко прието ребро множеството $A$ остава съгласувано.

Алгоритъмът никога не образува цикъл, понеже приема ребро само между различни компоненти. Тъй като графът е свързан, след разглеждане на всички ребра всички върхове принадлежат на една компонента. Тогава $A$ е свързан и ацикличен граф, тоест покриващо дърво. Понеже $A$ е съгласувано и само е покриващо дърво, то е минимално покриващо дърво.
\end{proof}

\section{Union--Find структури от данни}\label{s:7.5}

За ефикасната реализация на Крускал е необходима структура от данни, която поддържа разбиване на множество на непразни неприпокриващи се дялове.

\begin{defn}
Union--Find структура върху опорно множество $S$ поддържа следните операции:

\begin{itemize}
\item \verb|Make-Set(x)| създава едноелементния дял $\{x\}$;
\item \verb|Find-Set(x)| връща представител, наричан лидер, на дяла, съдържащ $x$;
\item \verb|Union(x,y)| слива двата дяла, съдържащи $x$ и $y$.
\end{itemize}
\end{defn}

В алгоритъма на Крускал всяка компонента на текущата гора е един дял. Така проверката дали реброто $\{u,v\}$ би образувало цикъл е проверката
\[
\operatorname{Find\mbox{-}Set}(u)
\neq
\operatorname{Find\mbox{-}Set}(v).
\]
Ако е вярна, реброто се приема и се изпълнява \verb|Union(u,v)|.

\subsection{Реализация с гора от коренови дървета}\label{s:7.5.1}

Разбиването може да се представи чрез ориентирана гора. Върховете на гората са елементите на $S$, а всяко дърво съответства на един дял. Коренът на дървото е лидерът на съответния дял. Нека $\pi[x]$ е родителят на $x$; за корен е изпълнено $\pi[x]=x$. Тази гора няма нищо общо с гората, която Крускал строи --- тя само записва кои върхове са в една компонента, а не кои ребра са приети.

Полето $\mathit{rank}[x]$ по-долу е горна граница за височината на дървото с корен $x$ и се използва само от евристиката \emph{union by rank}; при $\verb|Make-Set|$ то е $0$.

\begin{verbatim}
Make-Set(x)
    pi[x] <- x
    rank[x] <- 0

Find-Set(x)
    if x != pi[x] then
        return Find-Set(pi[x])
    else
        return x
\end{verbatim}

Без допълнителни правила операцията \verb|Union| може да образува дърво с линейна височина. Тогава \verb|Find-Set| би изисквала линейно време в най-лошия случай.

\subsection{Union by rank}\label{s:7.5.2}

\begin{defn}[Обединение по ранг]
При евристиката \emph{union by rank} коренът на дървото с по-малък ранг се закача към корена на дървото с по-голям ранг. При равни рангове единият корен се избира за нов корен и неговият ранг се увеличава с единица.
\end{defn}

\begin{verbatim}
Union(x,y)
    rx <- Find-Set(x)              // корените, а не самите x и y
    ry <- Find-Set(y)
    if rx = ry then
        return
    if rank[rx] < rank[ry] then
        pi[rx] <- ry
    else if rank[rx] > rank[ry] then
        pi[ry] <- rx
    else
        pi[ry] <- rx
        rank[rx] <- rank[rx] + 1
\end{verbatim}

\begin{lem}
Започвайки от едноелементни множества, след произволна последователност от операции \verb|Union| с union by rank височината на всяко дърво е $O(\log n)$, където $n$ е броят върхове в дървото.
\end{lem}

\begin{proof}
Поддържаме едновременно инвариантите, че височината на дървото е най-много ранга на корена и че дърво с ранг $h$ има поне $2^h$ върха. Те са верни за едноелементно дърво: височината и рангът са $0$, а броят върхове е $1$.

Ако корен с по-малък ранг се закачи към корен с по-голям ранг, новата височина е най-много по-големия ранг: височината на закаченото дърво плюс едно е най-много този ранг. Рангът на новия корен не се променя. При равни рангове $h$ новата височина е най-много $h+1$, а рангът на новия корен също става $h+1$.

За втория инвариант е достатъчно да разгледаме момента, в който рангът се увеличава от $h$ на $h+1$. Тогава се сливат две дървета с ранг $h$, всяко с поне $2^h$ върха. Новото дърво има поне
\[
2^h+2^h=2^{h+1}
\]
върха.

Следователно, ако дърво с $n$ върха има ранг $h$, то
\[
2^h\leq n,
\]
тоест
\[
h\leq \log_2 n.
\]
Понеже височината е най-много ранга, тя също е най-много $\log_2 n$. Затова \verb|Find-Set| работи за $O(\log n)$ време.
\end{proof}

\subsection{Компресия на пътя}\label{s:7.5.3}

\begin{defn}[Компресия на пътя]
Втората евристика е \emph{компресия на пътя}. При изпълнение на \verb|Find-Set(x)| всички върхове по намерения път към корена се насочват непосредствено към корена.
\end{defn}

\begin{verbatim}
Find-Set(x)
    if x != pi[x] then
        pi[x] <- Find-Set(pi[x])
    return pi[x]
\end{verbatim}

При едновременно използване на union by rank и компресия на пътя общата цена на последователност от \verb|Make-Set|, \verb|Find-Set| и \verb|Union| операции е
\[
O\bigl((n+m)\alpha(n)\bigr),
\]
където $\alpha$ е обратната функция на Акерман. Тя расте изключително бавно и за практически размери има много малка стойност.

\section{Сложност на алгоритъма на Крускал}\label{s:7.6}

Нека $n=|V|$ и $m=|E|$.

Инициализирането с $n$ операции \verb|Make-Set| отнема $O(n)$ време. Сортирането на ребрата е с време
\[
O(m\log m).
\]
При реализация с union by rank и компресия на пътя операциите на Union--Find за целия алгоритъм струват
\[
O\bigl((n+m)\alpha(n)\bigr).
\]
Следователно общата времева сложност е
\[
O\bigl(m\log m+(n+m)\alpha(n)\bigr).
\]

За прост неориентиран граф е изпълнено $m<n^2$, откъдето
\[
\log m=O(\log n).
\]
Сортирането доминира останалата работа, така че сложността на алгоритъма на Крускал обичайно се записва като
\[
O(m\log m)=O(m\log n).
\]

\begin{supp}[Бележка]
Ако се използва само union by rank, без компресия на пътя, операциите \verb|Find-Set| са $O(\log n)$, а алгоритъмът на Крускал отново има оценка $O(m\log n)$ поради сортирането. Компресията на пътя подобрява цената на операциите върху разбиването, макар че не променя доминиращата асимптотична оценка на Крускал при сравняващо сортиране.
\end{supp}

\section{Изпитен фокус}\label{s:7.7}

\begin{itemize}
\item Да се дефинират тегловен граф, покриващо дърво, тегло на дърво и минимално покриващо дърво.
\item Да се формулира задачата: неориентиран свързан граф с тегла на ребрата; търси се покриващо дърво с минимална сума от теглата.
\item Да се дефинират съгласувано множество, сигурно ребро, срез, съобразен срез и леко ребро.
\item Да се формулира и докаже теоремата за съгласуваното множество чрез добавяне на ребро, получаване на цикъл и размяна с друго ребро от същия срез.
\item Да се възпроизведе псевдокодът на Прим, значението на $key$, $\pi$ и приоритетната опашка, както и операциите \verb|Extract-Min| и \verb|Decrease-Key|, които той използва.
\item Да се обясни коректността на Прим чрез среза между вече избраните и неизбраните върхове и теоремата за съгласуваното множество.
\item Да се изведат сложностите на Прим: $O(m\log n)$ с двоична пирамида и $O(m+n\log n)$ с пирамида на Фибоначи.
\item Да се възпроизведе псевдокодът на Крускал и ролята на сортирането на ребрата.
\item Да се обясни коректността на Крускал: прието ребро свързва различни компоненти и е леко за подходящ съобразен срез.
\item Да се дефинират \verb|Make-Set|, \verb|Find-Set| и \verb|Union|, както и реализацията им чрез гора от коренови дървета.
\item Да се обяснят union by rank и компресията на пътя и приложението им в Крускал.
\item Да се изведе сложността на Крускал:
\[
O\bigl(m\log m+(n+m)\alpha(n)\bigr)=O(m\log n).
\]
\end{itemize}
